\subsection{Lexical decisions}

Keyword-rich query languages put most of the difficulty in the lexer rather than the
parser, and both grammars needed terminals shaped by word boundaries, lookaheads and
priorities. Table~\ref{tab:grammar-decisions} records them the same way
\S\ref{sec:datadog-inference} records the missing productions: each decision is
perturbed back to the form it would naively take, and a witness is re-parsed to show
what breaks. A row that still parsed would mean the decision was dead weight; none is.
The witnesses are minimal queries written for the table rather than corpus queries, which
an earlier version of this paragraph claimed they were; what the corpus supplies is the
count beside each one, of distinct corpus queries the real grammar parses and the naive
form does not. Three of the four decisions are load-bearing against the corpus as well as against
their witnesses. The fourth is not: the STAR terminal costs no corpus query at all,
because the construct that exercised its witness --- the keyless scope
\texttt{avg:x\{*-zone\}} --- is one \S\ref{sec:gaps} shows this paper's own normaliser
manufactured, so the decision is load-bearing against a string no dashboard wrote. An
earlier version of this paragraph asserted the universal and then named the colon
lookahead as the row where witness and corpus differ in kind; that row costs five
queries, and the row that differs is the one the sentence did not name. The colon
lookahead is still worth its own remark, because no corpus query names a metric
beginning with an aggregator word and what the corpus exhibits instead is a
\emph{function} whose name does, \texttt{count\_not\_null(...)}.

\begin{table}[htbp]\centering\scriptsize
\begin{tabular}{@{}llll@{}}
\toprule
Language & Decision & Witness & Naive form \\
\midrule
promql & keyword terminals carry \textbackslash{}b & \texttt{max\_over\_time(x[5m])} & UnexpectedToken \\
promql & FUNC needs a following paren & \texttt{up\{job="x"\}} & UnexpectedToken \\
datadog & aggregator looks ahead for the colon & \texttt{sum.of.things\{*\}} & UnexpectedToken \\
datadog & no separate STAR terminal & \texttt{avg:x\{*-zone\}} & UnexpectedToken \\
\bottomrule
\end{tabular}

\caption{Load-bearing lexical decisions. Each row perturbs one terminal to its naive
form and re-parses the witness, which the real grammar accepts. The witnesses are
minimal queries written for this table; the last column is how many distinct corpus
queries the naive form stops parsing.}
\label{tab:grammar-decisions}
\end{table}

% This file is itself deferred, so the hook inside it is unconditional: it is included
% exactly once either way, inline in \S\ref{sec:grammars} or in the appendix.
\begin{remark}[Priority before length, and the terminal it splits]\label{rem:lark-priority}
The last row is worth singling out because it contradicts a reasonable assumption. Lark
breaks a tie between two terminals that match at the same position by priority before
length, not longest-match-first: a high-priority \texttt{STAR} terminal for the match-all
scope \texttt{\{*\}} therefore splits the wildcard tag key \texttt{*-location} into
\texttt{STAR} followed by \texttt{TAG}. Folding the wildcard into the tag terminal and
letting the lowering layer recognise \texttt{*} is the fix, which makes match-all a
lexical property of a scope term rather than a production of its own.
\end{remark}

